% !TEX root = ../../hw01.tex
% Homework 01, Exercise 4. Shared solution.

% Transcribed from the lower left and upper right pages of IMG_8998.
\begin{proof}[Solution]
A map from $A$ to $B$ is a bijection if it is both one-to-one
and onto. This means that every element of $B$ is the image of
exactly one element of $A$.
\autoref{fig:hw01:bijections} contrasts bijections with a map
that fails both conditions.

\begin{figure}[h]
  \centering
  % !TEX root = ../../problems/hw01.tex
% Editable reconstruction of the finite-set mapping sketches in IMG_8998.
\begin{tikzpicture}[x=0.9cm,y=0.9cm,>=Stealth,font=\small]
  \foreach \panel/\captiontext in {0/{Bijection},3.8/{Bijection},7.6/{Not a bijection}} {
    \begin{scope}[shift={(\panel,0)}]
      \draw[black!55] (0,0) ellipse[x radius=0.4,y radius=1.05];
      \draw[black!55] (2.1,0) ellipse[x radius=0.4,y radius=1.05];
      \node at (0,1.35) {$A$};
      \node at (2.1,1.35) {$B$};
      \foreach \y in {-0.65,0,0.65} {
        \fill (0,\y) circle[radius=1.5pt];
        \fill (2.1,\y) circle[radius=1.5pt];
      }
      \node[font=\footnotesize] at (1.05,-1.4) {\captiontext};
    \end{scope}
  }
  \foreach \y in {-0.65,0,0.65}
    \draw[->,shorten <=3pt,shorten >=3pt] (0,\y) -- (2.1,\y);
  \draw[->,shorten <=3pt,shorten >=3pt] (3.8,0.65) -- (5.9,-0.65);
  \draw[->,shorten <=3pt,shorten >=3pt] (3.8,0) -- (5.9,0.65);
  \draw[->,shorten <=3pt,shorten >=3pt] (3.8,-0.65) -- (5.9,0);
  \draw[->,shorten <=3pt,shorten >=3pt] (7.6,0.65) -- (9.7,0.65);
  \draw[->,shorten <=3pt,shorten >=3pt] (7.6,0) -- (9.7,0.65);
  \draw[->,shorten <=3pt,shorten >=3pt] (7.6,-0.65) -- (9.7,-0.65);
\end{tikzpicture}

  \caption[Examples of bijections.]{The first two maps are bijections. The third map is neither
    one-to-one nor onto. Crossing arrows do not prevent a map from
    being bijective.}
  \label{fig:hw01:bijections}
\end{figure}

The tangent function is a bijection from
$(-\pi/2,\pi/2)$ onto $\R$. To use it on $(0,1)$, we first
multiply the input by $\pi$ and then subtract $\pi/2$.
These steps give the composition
\[
  (0,1)\xrightarrow{\,\times\pi\,}(0,\pi)
    \xrightarrow{\,-\pi/2\,}
    \left(-\frac\pi2,\frac\pi2\right)
    \xrightarrow{\,\tan\,}\R.
\]
We can therefore define
\[
  f(x)=\tan\left(\pi x-\frac\pi2\right),
    \qquad 0<x<1.
\]
For $0<x<1$, the angle $\pi x-\pi/2$ lies strictly between
$-\pi/2$ and $\pi/2$, so $f(x)$ is defined and real.
We verify bijectivity by giving its inverse. For $y\in\R$, put
\[
  g(y)=\frac12+\frac{\arctan y}{\pi}.
\]
Here $\arctan y$ is the unique angle in $(-\pi/2,\pi/2)$
whose tangent is $y$. Thus $0<g(y)<1$, and
\[
  f(g(y))=\tan(\arctan y)=y.
\]
This proves that $f$ is onto. Conversely, for $x\in(0,1)$,
the angle lies in the principal interval, so
\[
  g(f(x))=\frac12+
    \frac{\arctan(\tan(\pi x-\pi/2))}{\pi}=x.
\]
If $f(x_1)=f(x_2)$, applying $g$ gives $x_1=x_2$.
Therefore $f$ is also one-to-one and hence is a bijection.
\handoutonly{\clearpage}

\autoref{fig:hw01:tangent-overview} shows the central branch of
$y=\tan t$ in blue and the transformed graph alongside it.
The horizontal change $x=t/\pi+1/2$ compresses that branch by a
factor of $1/\pi$ and shifts it to the right by $1/2$.
Its domain becomes $(0,1)$, while its range remains all of $\R$.

\begin{figure*}[h]
  \centering
  % !TEX root = ../../problems/hw01.tex
% The tangent overview and its transformed central branch share one figure.
\begin{tikzpicture}[>=Stealth,font=\footnotesize]
  \definecolor{tangentblue}{HTML}{69B4E4}
  \definecolor{tangentwash}{HTML}{EAF6FD}

  \begin{scope}[x=0.62cm,y=0.42cm]
    \fill[tangentwash] (-1.5708,-4.6) rectangle (1.5708,4.6);
    \foreach \boundary in {-2.5,-1.5,1.5,2.5} {
      \draw[dashed,black!30] ({\boundary*pi},-4.6)
        -- ({\boundary*pi},4.6);
    }
    \foreach \boundary in {-0.5,0.5} {
      \draw[dashed,tangentblue] ({\boundary*pi},-4.6)
        -- ({\boundary*pi},4.6);
    }
    \draw[->,black!65] (-8.15,0) -- (8.3,0) node[right] {$t$};
    \draw[->,black!65] (0,-4.85) -- (0,5.05) node[above] {$y$};

    % Sample each branch separately, away from the vertical asymptotes.
    \foreach \period in {-2,-1,1,2} {
      \draw[thick,black!65,<->,domain=-1.347:1.347,samples=121,smooth]
        plot ({\x+\period*pi},{tan(deg(\x))});
    }
    \draw[line width=1.3pt,tangentblue,<->,
      domain=-1.347:1.347,samples=121,smooth]
      plot (\x,{tan(deg(\x))});

    \foreach \period/\ticktext in {-2/{-2\pi},-1/{-\pi},1/{\pi},2/{2\pi}} {
      \draw[black!65] ({\period*pi},0.12) -- ({\period*pi},-0.12);
      \node[below,fill=white,inner sep=2pt]
        at ({\period*pi},-0.14) {$\ticktext$};
    }
    \node[below right,inner sep=2pt] at (0,0) {$0$};
    \foreach \level in {-4,-2,2,4} {
      \draw[black!65] (-0.09,\level) -- (0.09,\level);
      \node[left,inner sep=2pt] at (-0.09,\level) {$\level$};
    }
    \foreach \boundary/\ticktext in {
      -2.5/{-5\pi/2},-1.5/{-3\pi/2},-0.5/{-\pi/2},
      0.5/{\pi/2},1.5/{3\pi/2},2.5/{5\pi/2}} {
      \node[below] at ({\boundary*pi},-4.85) {$\ticktext$};
    }
    \node at (0,-6.4) {$y=\tan t$};
  \end{scope}

  \begin{scope}[shift={(7.05cm,0)},x=2.75cm,y=0.42cm]
    \draw[dashed,black!45] (1,-4.6) -- (1,4.6);
    \draw[->,black!65] (-0.12,0) -- (1.18,0) node[right] {$x$};
    \draw[->,black!65] (0,-4.85) -- (0,5.05) node[above] {$y$};
    % The same parameter t places the selected branch at x=t/pi+1/2.
    \draw[thick,<->,domain=-1.347:1.347,samples=121,smooth]
      plot ({\x/pi+0.5},{tan(deg(\x))});
    \node[below left,fill=white,inner sep=1pt] at (0,0) {$0$};
    \node[below,fill=white,inner sep=2pt] at (1,0) {$1$};
    \node[below right,fill=white,inner sep=2pt] at (0.5,0) {$\tfrac12$};
    \node at (0.5,-6.4) {$y=\tan(\pi x-\pi/2)$};
  \end{scope}
\end{tikzpicture}

  \caption[Mapping an interval onto the real line.]{The left panel highlights the central tangent branch on
    $(-\pi/2,\pi/2)$. The right panel shows this branch after the
    change of input. The vertical asymptotes now occur at $x=0$
    and $x=1$. The horizontal axes use different scales.}
  \label{fig:hw01:tangent-overview}
\end{figure*}
\end{proof}
